\documentclass[conference]{IEEEtran}
\IEEEoverridecommandlockouts

\usepackage{cite}
\usepackage{amsmath,amssymb}
\usepackage{graphicx}
\usepackage{booktabs}
\usepackage{url}
\usepackage[hidelinks]{hyperref}

% Generated by rlcompare. Do not edit by hand.
\newcommand{\NumSeeds}{10}
\newcommand{\NumSeedsSweep}{5}
\newcommand{\NumTuningSeeds}{3}
\newcommand{\NumTrials}{12}
\newcommand{\StepBudget}{100k}
\newcommand{\EvalEvery}{2k}
\newcommand{\EvalEpisodes}{10}
\newcommand{\SolvedThreshold}{475}
\newcommand{\NumRunsMain}{30}
\newcommand{\NumRunsSweep}{75}
\newcommand{\NumRunsAblation}{20}
\newcommand{\NumRunsTuning}{108}
\newcommand{\NumRunsTotal}{233}
\newcommand{\TabFinal}{226}
\newcommand{\TabCiLow}{133}
\newcommand{\TabCiHigh}{390}
\newcommand{\TabTime}{16}
\newcommand{\TabUpdates}{100k}
\newcommand{\TabFailure}{10}
\newcommand{\TabSolve}{36k}
\newcommand{\TabLoc}{87}
\newcommand{\TabHyper}{7}
\newcommand{\DqnFinal}{486}
\newcommand{\DqnCiLow}{354}
\newcommand{\DqnCiHigh}{500}
\newcommand{\DqnTime}{180}
\newcommand{\DqnUpdates}{99k}
\newcommand{\DqnFailure}{0}
\newcommand{\DqnSolve}{55k}
\newcommand{\DqnLoc}{136}
\newcommand{\DqnHyper}{12}
\newcommand{\ReiFinal}{500}
\newcommand{\ReiCiLow}{493}
\newcommand{\ReiCiHigh}{500}
\newcommand{\ReiTime}{39}
\newcommand{\ReiUpdates}{470}
\newcommand{\ReiFailure}{0}
\newcommand{\ReiSolve}{18k}
\newcommand{\ReiLoc}{113}
\newcommand{\ReiHyper}{6}
\newcommand{\TabCoverage}{32}
\newcommand{\TabCells}{1\,024}
\newcommand{\TabEntries}{2\,048}
\newcommand{\DqnParams}{17\,410}
\newcommand{\ReiParams}{34\,691}
\newcommand{\AblNoBaselineFinal}{361}
\newcommand{\AblNoBaselineFailure}{20}
\newcommand{\AblNoBaselineSolve}{38k}
\newcommand{\AblNoTargetFinal}{9}
\newcommand{\AblNoTargetFailure}{100}
\newcommand{\AblNoTargetSolve}{---}
\newcommand{\TabSweepBest}{12x12x24x24}
\newcommand{\TabSweepBestValue}{401}
\newcommand{\TabSweepWorst}{3x3x6x6}
\newcommand{\TabSweepWorstValue}{131}
\newcommand{\TabSweepSpread}{270}
\newcommand{\DqnSweepBest}{500}
\newcommand{\DqnSweepBestValue}{418}
\newcommand{\DqnSweepWorst}{1}
\newcommand{\DqnSweepWorstValue}{9}
\newcommand{\DqnSweepSpread}{408}
\newcommand{\ReiSweepBest}{0.0003}
\newcommand{\ReiSweepBestValue}{500}
\newcommand{\ReiSweepWorst}{0.003}
\newcommand{\ReiSweepWorstValue}{459}
\newcommand{\ReiSweepSpread}{41}


\begin{document}

\title{A Comparative Study of Value-Based and Policy-Based\\
Reinforcement Learning on CartPole-v1}

% ---------------------------------------------------------------------------
% TODO(author): replace the institution, department and course code below.
% They are the only placeholders left in this document.
% ---------------------------------------------------------------------------
\newcommand{\Institution}{INSTITUTION NAME}
\newcommand{\Department}{Department of Computer Science}
\newcommand{\CourseCode}{COURSE CODE}

\author{\IEEEauthorblockN{Het Jethva}
\IEEEauthorblockA{\textit{\Department} \\
\textit{\Institution} \\
\CourseCode \\
hetjethva07@gmail.com}}

\maketitle

\begin{abstract}
% ABSTRACT
\end{abstract}

\begin{IEEEkeywords}
reinforcement learning, Q-learning, deep Q-network, policy gradient,
function approximation, empirical comparison
\end{IEEEkeywords}

\section{Introduction}

CartPole \cite{barto1983} is the standard first problem in reinforcement
learning. It is small, it is fast, and every general-purpose method solves it.
That last property is what makes it interesting for a comparison: because the
outcome is not in doubt, the question stops being \emph{which method works} and
becomes \emph{what each method costs to make it work}.

This paper compares three methods that sit at three levels of approximation:

\begin{itemize}
\item \textbf{Tabular Q-learning} \cite{watkins1992}, which represents the value
      function exactly over a discretised state space;
\item \textbf{DQN} \cite{mnih2015}, which replaces the table with a neural
      network and adds the two devices---a replay buffer and a target
      network---that make that replacement stable;
\item \textbf{REINFORCE with a learned baseline} \cite{williams1992}, which
      abandons the value function as the object of interest and optimises the
      policy directly.
\end{itemize}

The comparison is not a ranking. All three methods reach a working policy on
CartPole-v1, so a claim of the form ``method X is best'' would be both
uninformative and already known. The contribution here is a \emph{trade-off map}:
a measurement of what each method spends to arrive at that policy, along four
axes that are usually reported separately or not at all---environment samples,
computation, reliability across random seeds, and implementation effort.

Two design commitments make the resulting numbers comparable. First, every method
runs under a single shared training loop with an identical budget of
\StepBudget{} environment steps and an identical evaluation procedure, so no
method can gain an advantage from the harness. Second, hyperparameters are
selected by a random search with an equal trial budget per method, on seeds
disjoint from those used for the reported results. Without the second commitment,
a comparison paper measures how the author distributed their attention.

The contributions are:

\begin{enumerate}
\item A controlled comparison of the three methods on CartPole-v1 over
      \NumSeeds{} seeds each, under a shared budget and a fixed evaluation set.
\item A trade-off table quantifying sample cost, compute cost, reliability and
      implementation cost together, including the implementation-cost axis that
      empirical RL papers normally omit.
\item A sensitivity study for each method over the hyperparameter that carries
      its characteristic failure mode, and ablations of the two components that
      distinguish the deep methods from their textbook forms.
\item A complete, seed-reproducible implementation with the full result set.
\end{enumerate}

\section{Background}

\subsection{Problem setting}

Reinforcement learning problems are modelled as a Markov decision process
$(\mathcal{S}, \mathcal{A}, P, r, \gamma)$, where $P(s' \mid s, a)$ is the
transition distribution, $r(s,a)$ the reward, and $\gamma \in [0,1)$ a discount
factor. The agent seeks a policy $\pi(a \mid s)$ maximising the expected
discounted return $\mathbb{E}\!\left[\sum_{t} \gamma^{t} r_t\right]$
\cite{sutton2018}.

CartPole-v1 has a four-dimensional continuous observation---cart position, cart
velocity, pole angle and pole angular velocity---and two actions, pushing the
cart left or right. Every timestep yields a reward of $1$. An episode terminates
when the pole angle exceeds $12^{\circ}$ or the cart leaves the track, and is
truncated at 500 steps. The maximum achievable return is therefore 500, and the
conventional definition of solving the task is a mean return of at least
\SolvedThreshold{}.

\subsection{Tabular Q-learning}

Q-learning \cite{watkins1992} estimates the optimal action-value function by the
off-policy update
\begin{equation}
Q(s_t,a_t) \leftarrow Q(s_t,a_t) + \alpha
\left[ r_t + \gamma \max_{a'} Q(s_{t+1},a') - Q(s_t,a_t) \right],
\end{equation}
which converges to $Q^{*}$ under standard conditions when every state--action
pair is visited infinitely often and $Q$ is represented exactly.

That last requirement is the difficulty here. An exact representation needs a
finite state space, and CartPole's is continuous, so the observation must be
discretised into bins. Two of its four dimensions---the two velocities---are
unbounded, so they must additionally be clipped before they can be binned. The
resulting agent therefore does not solve CartPole; it solves a discretised
approximation of CartPole, and the gap between the two is a cost that the method
pays and the other two do not.

\subsection{Deep Q-networks}

DQN \cite{mnih2015} replaces the table with a parameterised function
$Q_\theta(s,a)$ and minimises the temporal-difference error against a bootstrap
target. Two devices stabilise what is otherwise a moving-target regression:

\begin{itemize}
\item a \emph{replay buffer}, which stores past transitions and samples training
      batches uniformly from them, breaking the temporal correlation that would
      otherwise violate the i.i.d.\ assumption of stochastic gradient descent;
\item a \emph{target network} $Q_{\theta^-}$, a periodically refreshed copy of
      the online network used to compute the bootstrap term, which stops the
      regression target from moving with every gradient step.
\end{itemize}

The loss on a sampled transition is
\begin{equation}
\mathcal{L}(\theta) = \ell_{\delta}\!\left(
  Q_\theta(s,a) -
  \left[ r + \gamma (1 - d) \max_{a'} Q_{\theta^-}(s',a') \right]
\right),
\end{equation}
where $\ell_{\delta}$ is the Huber loss and $d$ indicates a genuine termination.
Section~\ref{sec:extended} ablates the target network to show what it buys.

\subsection{REINFORCE with a baseline}

REINFORCE \cite{williams1992} optimises a parameterised policy $\pi_\theta$
directly, following the policy gradient
\begin{equation}
\nabla_\theta J(\theta) = \mathbb{E}_{\pi_\theta}
\left[ \nabla_\theta \log \pi_\theta(a_t \mid s_t)\,
\left( G_t - b(s_t) \right) \right],
\end{equation}
where $G_t$ is the Monte-Carlo return from timestep $t$ and $b(s_t)$ is a
baseline. Any function of the state alone leaves the gradient unbiased, so the
baseline is free to reduce variance; the standard choice, adopted here, is a
learned state-value estimate $\hat{V}_\phi(s)$ trained by regression on the
observed returns \cite{sutton2018}.

The baseline is not cosmetic. Because every CartPole reward is $+1$, the raw
return $G_t$ is large and strictly positive, so an unbaselined update increases
the log-probability of \emph{every} action taken in a successful episode,
including the poor ones---the gradient carries the correct signal only through
differences in magnitude. Subtracting $\hat{V}_\phi(s)$ centres the signal so
that only better-than-expected actions are reinforced. This is ablated in
Section~\ref{sec:extended}.

Note that returns are deliberately \emph{not} whitened. Dividing by the batch
standard deviation, common in tutorial implementations, rescales the gradient as
well as centring it; it is an adaptive step size rather than a baseline, and it
would confound the ablation.

\section{Experimental Design}
\label{sec:design}

The design follows the reproducibility recommendations of Henderson et al.\
\cite{henderson2018} and the statistical recommendations of Agarwal et al.\
\cite{agarwal2021}.

\subsection{Shared budget and a single training loop}

All three agents run under one training loop, differing only in what they do with
a transition. Each run receives exactly \StepBudget{} environment steps, and
learning curves are plotted against environment steps rather than episodes. The
distinction matters: a CartPole episode lasts around 20 steps for an untrained
policy and 500 for a good one, so an episode budget silently grants the better
method up to twenty-five times more data.

The loop bootstraps through a time-limit truncation but not through a genuine
termination. Treating the 500-step truncation as terminal would teach both
value-based methods that balancing the pole indefinitely is worth nothing.

\subsection{Evaluation protocol}

Training return is a poor measure of policy quality, because it includes the
cost of exploration and therefore penalises whichever method explores most. Every
\EvalEvery{} steps, training is paused and the agent is evaluated for
\EvalEpisodes{} episodes with exploration disabled: greedy action selection for
the two value-based methods, and the mode of the action distribution for
REINFORCE.

The evaluation episodes use a fixed set of environment seeds, identical for every
method, every seed and every evaluation point. The evaluation set is therefore
held out and constant, so a difference between two curves reflects the policies
rather than a difference in the states they happened to be tested on.

\subsection{Equal-budget hyperparameter search}

Each method receives \NumTrials{} random-search trials drawn from a discrete
space of its own hyperparameters, evaluated on \NumTuningSeeds{} tuning seeds and
scored by area under the evaluation curve---a criterion that rewards reaching a
high return and reaching it early. The winning configuration is then frozen and
run on \NumSeeds{} fresh evaluation seeds. Tuning seeds and evaluation seeds are
disjoint, so no reported result is tuned on its own test set. The full search
spaces and the selected values appear in Appendix~A.

\subsection{Statistics}

Reinforcement learning results across seeds are skewed and frequently bimodal: a
run either solves the task or collapses. The sample mean describes neither mode,
and one failed seed in ten moves it substantially. Following Agarwal et al.\
\cite{agarwal2021}, the point estimate is the \emph{interquartile mean} (IQM),
the mean of the middle 50\% of the seeds, and uncertainty is reported as a 95\%
percentile bootstrap interval resampled over seeds. Per-seed final performance is
also plotted directly, so that any bimodality remains visible rather than being
absorbed into an aggregate.

\section{Implementation}

The implementation is a single Python package. All three agents expose the same
four-operation interface---select an action, accept a transition, finish an
episode, report diagnostics---so the training loop, the evaluation procedure and
the logging are written once and shared. Every behavioural difference between the
methods lives in the update.

Both neural methods use a two-layer multilayer perceptron with 128 hidden units
per layer and ReLU activations, optimised with Adam. DQN uses this network for
$Q_\theta$ and a periodically synchronised copy for $Q_{\theta^-}$; REINFORCE
uses one for the policy logits and a second, separately optimised network for the
baseline $\hat{V}_\phi$.

One configuration object fully determines one run, and a hash of it names that
run's output files. This makes the run matrix resume-safe---an interrupted sweep
restarts without repeating completed work---and it means every curve in this
paper can be traced to the exact settings that produced it. Runs execute in a
process pool, one run per worker, each pinned to a single thread so that the
reported wall-clock times are not distorted by workers competing for cores.

Correctness is protected at the points where an error would be silent rather than
loud. A wrong sign in the advantage or an off-by-one in the discounted return
does not raise an exception; it produces a plausible learning curve and a wrong
paper. The discounted-return computation, the state discretiser's clipping
behaviour, the replay buffer's contents under overwrite, and the aggregation
statistics are therefore checked against hand-computed values, and every method
is checked to reproduce its curve exactly from a fixed seed.

Every figure and table in this paper is generated from the raw result files, and
every number quoted in the text is a macro produced by the same stage, so the
prose cannot drift out of step with the data.

\section{Results}
% RESULTS

\section{Extended Analysis}
% EXTENDED

\section{Discussion}
% DISCUSSION

\section{Threats to Validity}
% THREATS

\section{Conclusion}
% CONCLUSION

\bibliographystyle{IEEEtran}
\bibliography{references}

\appendices
\section{Hyperparameter Settings}
\begin{table}[t]
\centering
\caption{stub}
\label{tab:hyperparameters}
\begin{tabular}{llc}
\toprule
A & B & C \\
\midrule
x & y & z \\
\bottomrule
\end{tabular}
\end{table}


\end{document}
